%!TeX root=../emmatop.tex
\chapter[Chapter \thechapter]{}
\lettrine[lraise=0.3]{I}{f} Emma had still, at intervals, an anxious feeling for Harriet, a momentary doubt of its being possible for her to be really cured of her attachment to Mr Knightley, and really able to accept another man from unbiased inclination, it was not long that she had to suffer from the recurrence of any such uncertainty. A very few days brought the party from London, and she had no sooner an opportunity of being one hour alone with Harriet, than she became perfectly satisfied—unaccountable as it was!—that Robert Martin had thoroughly supplanted Mr Knightley, and was now forming all her views of happiness.

Harriet was a little distressed—did look a little foolish at first: but having once owned that she had been presumptuous and silly, and self-deceived, before, her pain and confusion seemed to die away with the words, and leave her without a care for the past, and with the fullest exultation in the present and future; for, as to her friend's approbation, Emma had instantly removed every fear of that nature, by meeting her with the most unqualified congratulations.—Harriet was most happy to give every particular of the evening at Astley's, and the dinner the next day; she could dwell on it all with the utmost delight. But what did such particulars explain?—The fact was, as Emma could now acknowledge, that Harriet had always liked Robert Martin; and that his continuing to love her had been irresistible.—Beyond this, it must ever be unintelligible to Emma.

The event, however, was most joyful; and every day was giving her fresh reason for thinking so.—Harriet's parentage became known. She proved to be the daughter of a tradesman, rich enough to afford her the comfortable maintenance which had ever been hers, and decent enough to have always wished for concealment.—Such was the blood of gentility which Emma had formerly been so ready to vouch for!—It was likely to be as untainted, perhaps, as the blood of many a gentleman: but what a connexion had she been preparing for Mr Knightley—or for the Churchills—or even for Mr Elton!—The stain of illegitimacy, unbleached by nobility or wealth, would have been a stain indeed.

No objection was raised on the father's side; the young man was treated liberally; it was all as it should be: and as Emma became acquainted with Robert Martin, who was now introduced at Hartfield, she fully acknowledged in him all the appearance of sense and worth which could bid fairest for her little friend. She had no doubt of Harriet's happiness with any good-tempered man; but with him, and in the home he offered, there would be the hope of more, of security, stability, and improvement. She would be placed in the midst of those who loved her, and who had better sense than herself; retired enough for safety, and occupied enough for cheerfulness. She would be never led into temptation, nor left for it to find her out. She would be respectable and happy; and Emma admitted her to be the luckiest creature in the world, to have created so steady and persevering an affection in such a man;—or, if not quite the luckiest, to yield only to herself.

Harriet, necessarily drawn away by her engagements with the Martins, was less and less at Hartfield; which was not to be regretted.—The intimacy between her and Emma must sink; their friendship must change into a calmer sort of goodwill; and, fortunately, what ought to be, and must be, seemed already beginning, and in the most gradual, natural manner.

Before the end of September, Emma attended Harriet to church, and saw her hand bestowed on Robert Martin with so complete a satisfaction, as no remembrances, even connected with Mr Elton as he stood before them, could impair.—Perhaps, indeed, at that time she scarcely saw Mr Elton, but as the clergyman whose blessing at the altar might next fall on herself.—Robert Martin and Harriet Smith, the latest couple engaged of the three, were the first to be married.

Jane Fairfax had already quitted Highbury, and was restored to the comforts of her beloved home with the Campbells.—The Mr Churchills were also in town; and they were only waiting for November.

The intermediate month was the one fixed on, as far as they dared, by Emma and Mr Knightley.—They had determined that their marriage ought to be concluded while John and Isabella were still at Hartfield, to allow them the fortnight's absence in a tour to the seaside, which was the plan.—John and Isabella, and every other friend, were agreed in approving it. But Mr Woodhouse—how was Mr Woodhouse to be induced to consent?—he, who had never yet alluded to their marriage but as a distant event.

When first sounded on the subject, he was so miserable, that they were almost hopeless.—A second allusion, indeed, gave less pain.—He began to think it was to be, and that he could not prevent it—a very promising step of the mind on its way to resignation. Still, however, he was not happy. Nay, he appeared so much otherwise, that his daughter's courage failed. She could not bear to see him suffering, to know him fancying himself neglected; and though her understanding almost acquiesced in the assurance of both the Mr Knightleys, that when once the event were over, his distress would be soon over too, she hesitated—she could not proceed.

In this state of suspense they were befriended, not by any sudden illumination of Mr Woodhouse's mind, or any wonderful change of his nervous system, but by the operation of the same system in another way.—Mrs Weston's poultry-house was robbed one night of all her turkeys—evidently by the ingenuity of man. Other poultry-yards in the neighbourhood also suffered.—Pilfering was housebreaking to Mr Woodhouse's fears.—He was very uneasy; and but for the sense of his son-in-law's protection, would have been under wretched alarm every night of his life. The strength, resolution, and presence of mind of the Mr Knightleys, commanded his fullest dependence. While either of them protected him and his, Hartfield was safe.—But Mr John Knightley must be in London again by the end of the first week in November.

The result of this distress was, that, with a much more voluntary, cheerful consent than his daughter had ever presumed to hope for at the moment, she was able to fix her wedding-day—and Mr Elton was called on, within a month from the marriage of Mr and Mrs Robert Martin, to join the hands of Mr Knightley and Miss Woodhouse.

The wedding was very much like other weddings, where the parties have no taste for finery or parade; and Mrs Elton, from the particulars detailed by her husband, thought it all extremely shabby, and very inferior to her own.—<Very little white satin, very few lace veils; a most pitiful business!—Selina would stare when she heard of it.>—But, in spite of these deficiencies, the wishes, the hopes, the confidence, the predictions of the small band of true friends who witnessed the ceremony, were fully answered in the perfect happiness of the union.

\vfill
\begin{figure}[bh]
\centering
\includegraphics[width=.5\linewidth]{cameo}
\end{figure}
